\section{模型检验与灵敏度分析}
\subsection{硬约束、复算与双运行}
solve v13 的约束审计显示：任务唯一性、到达、截止、实时开工、时延、GPU、IT 与设施功率均通过；Q3 和 Q4 的能源平衡、SOC、功率、互斥及终端状态最大绝对残差不超过 $1.137\times10^{-13}$。这些检验说明结果在给定容差与题面参数内自洽，但不提供最优性证书。

两次隔离运行使用相同输入、同一 code v16 与固定种子，自然耗时分别为 754.37 s 和 793.71 s，输出文件数均为 53。\texttt{results.json} 的 SHA-256 同为 \texttt{377f0fc8...759956}；Q2、Q3、Q4 结构化语义及去除 run ID、真实耗时后的运行合同完全一致。仅三份包含实测耗时的文件字节不同。

\subsection{问题二参数灵敏度}
以平衡方案成本 $-421378290.28$ 元为基准，对购电价格整体缩放。系数为 0.8、0.9、1.1、1.2 时，成本分别为 $-421379709.76$、$-421379000.02$、$-421377580.54$、$-421376870.80$ 元，相对变化绝对值不超过 0.000337\%。图~\ref{fig:price} 表明当前净成本对购电价格整体缩放不敏感，原因是购电支出相对售电收入很小。
\begin{figure}[H]\centering
\includegraphics[width=.80\textwidth]{sensitivity_electricity_price.png}
\caption{问题二净成本的购电价格灵敏度}\label{fig:price}
\end{figure}

新能源水平系数为 0.8、0.9、1.1、1.2 时，成本分别为 $-3.046617$、$-3.806991$、$-4.462847$、$-4.576383$ 亿元，相对基准变化为 27.6988\%、9.6539\%、$-5.9107$\%、$-8.6051$\%。图~\ref{fig:renewable} 呈现明显非线性：新能源下降时购电缺口扩大，而上升时外送上限与弃电共同限制边际收益。
\begin{figure}[H]\centering
\includegraphics[width=.80\textwidth]{sensitivity_renewable_level.png}
\caption{问题二净成本的新能源水平灵敏度}\label{fig:renewable}
\end{figure}

\subsection{终端 SOC 与高循环策略解释}
15、21、31 个 SOC 状态下，统一所选成本策略的成本分别为 $-4.583249$、$-4.587522$、$-4.590387$ 亿元；21 相对 31 状态的成本误差为 0.0624\%，新能源利用率误差为 0.0547\%，且三组约束残差均不超过 $1.137\times10^{-13}$。对应等效循环为 247.737、240.186、231.961。题面未给退化成本与循环上限，故这些高吞吐轨迹只表示理想设备情景可行。

Q4 三个新能源波动种子的峰值购电为 131.0484--292.4869 MW、碳排放为 333.5577--636.6865 tCO$_2$，显示单一随机轨迹不足以概括风险；当前三种子只提供方向性稳健性证据。
